% ===================================================================
%  HOTO Na 10 — Teku. --- Tashar jiragen ruwa.
%
%  Traduction, faite en 2026, en haoussa d'aujourd'hui, sur l'ido de
%  texto/io. Regles de la colonne : voir 10-hoto-01.tex. Une seule
%  numerotation.
%
%  « jirgi » EST LE SIXIEME GRAND MOT PRODUCTIF DE CETTE COLONNE, ET
%  IL PARTAGE LE MONDE EN TROIS ELEMENTS. Un jirgi est n'importe quel
%  grand porteur, et c'est le milieu qui le distingue : jirgin ruwa le
%  bateau — le porteur d'eau —, jirgin ƙasa le train — celui de terre,
%  qu'on verra au tableau 12 —, jirgin sama l'avion, celui du ciel.
%  Tout ce tableau-ci est ecrit avec le premier des trois : jirgin
%  ruwa mai nutsewa le sous-marin (10), jirgin ruwa mai sulke le
%  cuirasse (9), jiragen kamun kifi les barques de peche (55), jirgin
%  jan wasu le remorqueur (16). La ou l'ido a quinze racines
%  differentes, le haoussa a un mot et des complements.
%
%  ET « keke » ATTEINT SIX : les vapeurs a aubes (31) sont « jiragen
%  ruwa masu keke ». Bicyclette, voiture attelee, rouet, roue de
%  moulin, funiculaire, roue a aubes. Il n'y a toujours rien de commun
%  entre ces six objets qu'un axe.
%
%  DEUX CHAMEAUX, ET C'EST LA PLUS BELLE PAIRE DU RELEVE. La girafe du
%  tableau 3 etait « raƙumin dawa », le chameau de la brousse ; la
%  vague (72) est « raƙumin ruwa », le chameau de l'eau. Une langue du
%  desert mesure sur la seule bete qu'elle connaisse bien tout ce qui
%  est trop haut et tout ce qui roule sans fin. Aucune autre colonne
%  du releve ne nomme deux choses aussi eloignees par le meme animal.
%
%  HUITIEME REFUS DE LA COLONNE : LES BARQUES (46, 49, 54, 55).
%  « kwale-kwale » est la pirogue du Niger et de la Benoue, creusee
%  dans un seul tronc, et c'est le mot qu'un Haoussa emploierait sans
%  y penser devant n'importe quelle petite embarcation. La planche
%  grave des yoles bordees, membrees et pontees, un canot a voile
%  avec vergues et haubans, des barques de peche a quille. Ce ne sont
%  pas des troncs creuses. On ecrit « ƙananan jiragen ruwa ».
%
%  ET TOUTE LA MARINE PARLE ANGLAIS, TANDIS QUE LA MER EST HAOUSSA.
%  D'un cote laftana, kyaftin, admiral, kwamanda, kwastan, tiram,
%  anka, fanka, turfido : pas un grade, pas une machine, pas un
%  bureau qui ne soit venu avec son nom. De l'autre teku la mer, ruwa
%  l'eau, raƙumin ruwa la vague, tsibiri l'ile, gaci la falaise,
%  guguwa la tempete, bakin teku la plage. C'est la loi du tableau 3,
%  verifiee une quatrieme fois : l'INSTITUTION arrive avec ses mots,
%  la CHOSE garde les siens.
%
%  ET LE CUIRASSE (9) SE DIT « jirgin ruwa mai sulke ». Le sulke est
%  l'armure matelassee que la cavalerie de Kano portait encore au
%  XXe siecle, plusieurs epaisseurs de coton piquees, assez epaisses
%  pour arreter une fleche. Le mot passe du cheval au vapeur sans
%  qu'on ait rien a lui ajouter.
%
%  ENFIN « fitila » PARAIT POUR LA SIXIEME FOIS : le phare (93) est
%  « fitilar teku », la lampe de la mer. Lampe a huile de l'atelier
%  de dessin au tableau 1, lampe a petrole du salon et lampe
%  electrique au tableau 6, lanterne de la cuisine, veilleuse de la
%  chambre — et une tour sur un rocher.
% ===================================================================

%%K t10-apar-1 apar
\VUsaut{4.0mm}
\VUtitre{15.0pt}{60}{HOTO Na 10}
\VUsaut{3.0mm}
\VUpk{11.4pt}{40}{Teku. --- Tashar jiragen ruwa.}
\VUsaut{3.0mm}

%%K t10-01-1 p
1. --- A lokacin rani, alhali waɗansu suna neman natsuwa da sanyi a
kan dutse, waɗansu sun fi son zuwa gaɓar teku. Haka muka yi bara,
domin muna son \VUgras{Babban Teku} \textsuperscript{(1)} ƙwarai.

%%K t10-02-1 p
2. --- Game da ni, ina jin daɗi ƙwarai ina kallon wannan faɗin ruwa
mai girma wanda yake miƙewa har \VUgras{iyakar ido}
\textsuperscript{(2)}, ina kuma ganin \VUgras{manyan jiragen ruwa masu
tururi} \textsuperscript{(3)} suna tashi domin doguwar \VUgras{tafiya,}
jiragen da matafiya masu zuwa Amirka suke hawa.

%%K t10-03-1 p
3. --- Shi ya sa mahaifina, wanda ya san abin da nake so, kullum yakan
zaɓi, domin mu yi hutu, wani bakin teku mai natsuwa ƙwarai, wanda yake
kusa da ɗaya daga cikin manyan \VUgras{tashoshin jiragen ruwa na
yaƙinmu.}

%%K t10-03-2 p
A zamanmu na ƙarshe, \VUgras{rundunar jiragen ruwa} ta zo ta kafa anka
a bayan babbar \VUgras{katangar ruwa} \textsuperscript{(4)} wadda take rufe
kuma take kare \VUgras{tashar jiragen ruwa ta yaƙi}
\textsuperscript{(5)}, na kuma yi nasarar yaba wa \VUgras{jiragen ruwan
yaƙi} iri iri kamar yadda nake so, wato ƙaramin \VUgras{jirgin ruwa mai
nutsewa} \textsuperscript{(10)} wanda yake nutsewa ya ɓace a ƙarƙashin
ruwa da kuma babban \VUgras{jirgin ruwa mai sulke} \textsuperscript{(9)}
wanda, har yanzu, kusan harin \VUgras{jirgin ruwan turfido}
\textsuperscript{(8)} kaɗai yake tsoro.

%%K t10-tit-2 sub
\VUsaut{3.0mm}
\VUpk{10.2pt}{40}{Ƙananan jiragen ruwa.}
\VUsaut{3.0mm}

%%K t10-04-1 p
4. --- Wannan ya riga ya san yadda zai sami wurin da ya raunana na
sulken ƙarfe mai kauri \VUgras{cikin gwaninta,} domin ya kafa a can
\VUgras{turfido} mai haɗari, wanda fashewarsa take sa jirgin ya nutse,
amma duk da haka ba a jin tsoronsa kamar jirgin ruwa mai nutsewa,
domin jiragen hana turfido suna iya korarsa cikin sauƙi.

%%K t10-05-1 p
5. --- Na kuma iya ganin \VUgras{jiragen sintiri} \textsuperscript{(6)}
masu sauri da aka yi wa sulke, waɗanda kusan ba a iya ji musu ciwo
kamar jirgin ruwa mai sulke na gaskiya, da \VUgras{jiragen ɗaukar kaya}
\textsuperscript{(12)} da yawa, da \VUgras{jiragen igwa}
\textsuperscript{(7)} uku ko huɗu, a ƙarshe kuma \VUgras{faragat}
\textsuperscript{(11)} ɗaya. \VUgras{Kallon wannan rundunar jiragen ruwa,
idan take motsi domin ta ɗauke anka, shi ne mafi ban sha’awa.}

%%K t10-06-1 p
6. --- Kai, bambancin gini nawa ne tsakanin waɗannan manyan tuddan
ƙarfe da kyawawan \VUgras{jiragen mastu uku} ko \VUgras{jiragen ruwa
masu zanen iska} \textsuperscript{(13)} masu kwarjini waɗanda har
yanzu ake amfani da su a rundunar kasuwanci, waɗanda na gan su suna
shiga tashar, da dukan zanen iskarsu, lokacin da nake yawo a kan
\VUgras{gaɓar tashar} \textsuperscript{(14)}. Hakika waɗannan sun rasa
amfaninsu da yawa \VUgras{lokacin} da \VUgras{iska} ta zama ta gaba.
Dole su saukar da dukan zanen iska domin su sake shiga tashar,
\VUgras{jirgin jan wasu} \textsuperscript{(16)} kuwa ya zama dole domin ya
ɗauke su ya kai su, kamar \VUgras{jiragen kaya}
\textsuperscript{(17)} kawai, zuwa tashar ruwa.

%%K t10-tit-3 sub
\VUsaut{3.0mm}
\VUpk{10.2pt}{40}{Tashar jiragen ruwa ta kasuwanci.}
\VUsaut{3.0mm}

%%K t10-07-1 p
7. --- Abin da yake ba da sha’awa a tashar da nake magana a kanta shi
ne cewa ba tashar yaƙi kaɗai take ɗauke da ita ba, wadda aka shirya da
hannu da ayyuka masu girma, sai dai kuma \VUgras{tashar kasuwanci} mai
ban mamaki wadda take a \VUgras{bakin} babban kogi.

%%K t10-08-1 p
8. --- Harshen ƙasar da yake raba tashoshi biyun an yi masa kagara ana
kuma kare shi da jerin \VUgras{wuraren igwa} da \VUgras{hasumiyoyin
yaƙi} waɗanda suke miƙewa a kan teku. Ana buƙatar izini na musamman
domin a yi yawo a kan \VUgras{katangar yaƙi} \textsuperscript{(15)}. Na
sami shi cikin sauƙi, ta wurin baffana, wanda yake injiniyan
\VUgras{wuraren gina jiragen ruwa} \textsuperscript{(18)}, na kuma je na
ziyarci jiragen da ake gina a ƙarƙashin manyan rumfuna.

%%K t10-09-1 p
9. --- Har na halarci saukar da jirgin ruwa mai sulke a ruwa, ina kuma
tuna lokacin motsin rai da dukan taron suka ji sa’ad da katon abin nan,
da aka bar shi ga kansa, ya fara zamiya a kan katakon da yake riƙe da
shi ya kuma zo ya yi iyo a kan ruwan kogin.

%%K t10-10-1 p
10. --- Domin a je wuraren gina jiragen ruwa, akan ƙetare \VUgras{filin
atisaye} \textsuperscript{(19)} inda sojojin sansanin suke zuwa horo kowace
rana, akan wuce a kan \VUgras{ƙaramar gada} \textsuperscript{(20)} ta
ƙarfe, wadda take haɗa filin faretin da \VUgras{ma’ajin makamai}
\textsuperscript{(21)}.

%%K t10-tit-4 sub
\VUsaut{3.0mm}
\VUpk{10.2pt}{40}{Ma’ajin makamai da wuraren gyara.}
\VUsaut{3.0mm}

%%K t10-11-1 p
11. --- Tare da baffana, na ziyarci wannan babban gini cike da
\VUgras{igwaye} \textsuperscript{(22)}, da \VUgras{harsasan igwa}
\textsuperscript{(23)} da \VUgras{bama-bamai.} Na kuma ga \VUgras{wurin
narkar da ƙarfe} \textsuperscript{(24)} inda ake yin \VUgras{manyan
tukwanen tururi} \textsuperscript{(25)} na jiragen ruwa masu tururi.

%%K t10-12-1 p
12. --- A kusa ƙwarai, akwai \VUgras{wuraren gyaran jiragen ruwa}
\textsuperscript{(26)} da \VUgras{wurin gyara mai iyo}
\textsuperscript{(27)} mai ban mamaki ƙwarai, inda na iya gani a kusa
ƙwarai, a kan jirgin da za a gyara, \VUgras{fankar}
\textsuperscript{(28)}, da \VUgras{ƙasan jirgin} \textsuperscript{(29)} da
\VUgras{linzamin} \textsuperscript{(30)} jirgin.

%%K t10-13-1 p
13. --- Tashar kasuwanci tana da daraja a yi nazarinta kamar tashar
yaƙi, na kuma yi sa’o’i da yawa ina kallo a kan tashoshin ruwa inda aka
ɗaure \VUgras{jiragen ruwa masu keke} \textsuperscript{(31)} waɗanda suke
hawan kogin, ina kuma kallon \VUgras{manyan injinan ɗagawa}
\textsuperscript{(32)} suna aiki, waɗanda suke ɗaga manyan kaya kamar
gashin tsuntsu.

%%K t10-tit-5 sub
\VUsaut{4.0mm}
\VUpk{10.2pt}{40}{Matuƙin jirgin ruwa.}
\VUsaut{3.0mm}

%%K t10-14-1 p
14. --- Na yaba wa \VUgras{jiragen tsallaka teku} \textsuperscript{(34)}
suna fita daga mashigin, da \VUgras{tutocinsu} \textsuperscript{(35)} cikin
iska, \VUgras{matuƙin jirgin ruwa} \textsuperscript{(36)} gwani yake
jagorantarsu, shi wanda ya iya bin \VUgras{hanyar ruwa} da aka yi mata
alama da \VUgras{alamun ruwa} \textsuperscript{(40)} da sandunan nuni cikin
gwaninta, yana guje wa \VUgras{jiragen kaya masu zanen iska}
\textsuperscript{(41)} waɗanda ake tuƙa su da wahala da zanen iskarsu
guda ɗaya mai kusurwa huɗu. Na rarrabe a \VUgras{gaban jirgin}
\textsuperscript{(37)} — hancin jirgi — \VUgras{ɗakunan}
\textsuperscript{(39)} na aji na biyu, kuma a \VUgras{bayan jirgin}
\textsuperscript{(38)} — wutsiyar jirgi — falo domin matafiya na aji na
farko.

%%K t10-15-1 p
15. --- Wani lokaci na kuma halarci ɗora kaya a kan \VUgras{jirgin
kaya} \textsuperscript{(42)} inda aka tara kayayyakin \VUgras{ma’aji}
\textsuperscript{(43)} da na \VUgras{tashar teku} \textsuperscript{(44)}
yanzu, waɗanda jiragen ƙasa masu yawa na \VUgras{layin dogo na tashar}
\textsuperscript{(45)} suka kawo.

%%K t10-tit-6 sub
\VUsaut{3.0mm}
\VUpk{10.2pt}{40}{Tuƙi domin jin daɗi.}
\VUsaut{3.0mm}

%%K t10-16-1 p
16. --- Sau da yawa na yi tuƙi a \VUgras{tashar jiragen ƙanana}
\textsuperscript{(53)}, inda \VUgras{ƙananan jiragen ruwa}
\textsuperscript{(46)} suke zuwa su huta, kuma jin daɗi ne a gare ni in
motsa \VUgras{sandar tuƙi} \textsuperscript{(47)}. Na hau \VUgras{benen
jirgi} \textsuperscript{(48)} na \VUgras{ƙaramin jirgi mai zanen iska}
\textsuperscript{(49)}, na hau, ta wurin \VUgras{manyan igiyoyi}
\textsuperscript{(52)}, a kan \VUgras{mastu} \textsuperscript{(51)} zuwa
\VUgras{sandunan zane} \textsuperscript{(50)}, sa’an nan na sake yawona
\VUgras{da ƙaramin jirgi} \textsuperscript{(54)} tsakanin \VUgras{jiragen
kamun kifi} \textsuperscript{(55)} masu nauyi, inda \VUgras{taruna}
\textsuperscript{(56)} suke bushewa.

%%K t10-17-1 p
17. --- A kan \VUgras{tashar ruwa} \textsuperscript{(58)},
\VUgras{kayayyaki} \textsuperscript{(59)} iri iri sun wuce a gabana,
waɗanda suke cikin \VUgras{akwatuna} \textsuperscript{(60)} ko cikin
\VUgras{manyan ɗamummuka} \textsuperscript{(61)}. Sun yi \VUgras{kayan}
\textsuperscript{(63)} manyan jiragen kaya, \VUgras{injinan ɗagawa}
\textsuperscript{(62)} masu ƙarfi kuwa suka ɗauke su suka ajiye su a kan
\VUgras{manyan motoci} \textsuperscript{(64)} alhali \VUgras{masu jigilar
kaya} suka bayyana su suka kuma biya haraji a \VUgras{ofishin kwastan}
\textsuperscript{(65)}.

%%K t10-18-1 p
18. --- A ƙarshe, lokacin da na iso \VUgras{gada mai juyawa}
\textsuperscript{(66)} ko \VUgras{ƙofar ruwa} \textsuperscript{(69)}, ina
wucewa gaban \VUgras{injin tsintar ƙasa} \textsuperscript{(68)} wanda yake
tsaftace \VUgras{magudanar ruwa} \textsuperscript{(67)}, sai na sauka a
gaɓar tashar kusa da \VUgras{alamar jiragen ruwa}
\textsuperscript{(70)}.

%%K t10-19-1 p
19. --- A ɗaya gefen wannan gaɓa ne \VUgras{ƙananan jiragen sojoji}
\textsuperscript{(71)} suke zuwa su sauka, waɗanda suke kai jami’an
rundunar ƙasa. Abin sha’awa ne mutum ya ga matuƙan jirgin ruwa suna
ɗaga sandunan tuƙinsu cikin kari, alhali jirgin da \VUgras{raƙumin
ruwa} \textsuperscript{(72)} yake ɗauke da shi yake sauka a matakalar
wurin sauka cikin sauƙi. A wannan wuri, mutum yakan iya lura da
\VUgras{gudun teku} da motsin \VUgras{hawan teku} da \VUgras{jan bayansa}
a kan \VUgras{bakin teku} \textsuperscript{(73)} da kyau.

%%K t10-tit-7 sub
\VUsaut{3.0mm}
\VUpk{10.2pt}{40}{Gidan ƙungiya.}
\VUsaut{3.0mm}

%%K t10-20-1 p
20. --- Domin in komo gida, na yi amfani da \VUgras{tiram na lantarki}
\textsuperscript{(74)} wanda \VUgras{tashar sa} \textsuperscript{(75)} take
kusa da \VUgras{gungumen} da aka ajiye gaban ƙungiyar jami’ai.

%%K t10-20-2 p
Daga kan \VUgras{barandar sama} \textsuperscript{(76)} ta ƙungiyar, wadda
aka yi wa ado da \VUgras{ankoki} \textsuperscript{(77)} da igwaye da
harsasan igwa, sau da yawa na ga taron jami’an ruwa masu kwarjini :
\VUgras{laftanoni} \textsuperscript{(78)}, da takobinsu, da
\VUgras{kyaftinonin igwa} \textsuperscript{(79)} da \VUgras{na sojojin
ƙasa} \textsuperscript{(80)} da \VUgras{takobinsu mai lanƙwasa.} A
bayansu akwai \VUgras{matuƙin jirgin ruwa} \textsuperscript{(81)} wanda
yake kawo musu \VUgras{madubin gani mai nisa,} idan suna son su
rarrabe abin da yake faruwa can nesa.

%%K t10-21-1 p
21. --- Na kuma ga matasan \VUgras{ƙananan laftanoni}
\textsuperscript{(82)} suna karɓar umarnin \VUgras{kwamandansu}
\textsuperscript{(83)} ko na \VUgras{admiral} \textsuperscript{(84)},
alhali \VUgras{sajan jirgin ruwa} \textsuperscript{(85)} yake jira domin
ya kai su ga jirgin.

%%K t10-22-1 p
22. --- Daga wannan baranda, kallon yana da kyau ƙwarai : mutum yana
hango dukan bakin teku da \VUgras{tantinsa} \textsuperscript{(86)} da
\VUgras{ƙananan ɗakunansa} \textsuperscript{(87)}, mutum kuma yana gani,
a lokacin wanka, \VUgras{masu wanka} \textsuperscript{(88)} suna yin wasa
da farin ciki gaban \VUgras{gidan caca} \textsuperscript{(89)}.

%%K t10-23-1 p
23. --- A ƙarshen bakin teku, gaɓar tana yin ƙaramin \VUgras{hancin
ƙasa} \textsuperscript{(91)} \VUgras{mai duwatsu,} inda \VUgras{manyan
raƙuman ruwa} \textsuperscript{(92)} suke zuwa su karye kuma inda aka gina
\VUgras{fitilar teku} \textsuperscript{(93)}, wadda take nuna hanya ga
matuƙan jirgi da dare. Wannan hancin ƙasa yana haɗuwa da ƙasar ta
\VUgras{matsattsiyar hanyar ƙasa} kaɗai \textsuperscript{(90)}.

%%K t10-tit-8 sub
\VUsaut{3.0mm}
\VUpk{10.2pt}{40}{Kallon gaɓar teku gaba ɗaya.}
\VUsaut{3.0mm}

%%K t10-24-1 p
24. --- \VUgras{Gaci mai tsawo} \textsuperscript{(94)} yana miƙewa, teku
ya tsaga shi, teku wanda yake zana a cikinsa jerin \VUgras{mashiga}
\textsuperscript{(95)} da \VUgras{hancin ƙasa} yadda ya ga dama, yana sa
shi ya zama mai ban sha’awa ƙwarai.

%%K t10-24-2 p
A kan \VUgras{tudun fili} \textsuperscript{(96)}, wanda yake hangen
\VUgras{mashigin teku} \textsuperscript{(98)}, akwai \VUgras{kagara}
\textsuperscript{(97)} mai muhimmanci ƙwarai da waɗansu « gidajen
hutu ».

%%K t10-25-1 p
25. --- Wannan mashigi yana raba da nahiya \VUgras{tsibiri}
\textsuperscript{(99)} mai haɗari ƙwarai, wanda \VUgras{duwatsun ruwa}
\textsuperscript{(100)} da wuraren karyewar raƙuma suka kewaye shi, inda
ƙananan jiragen ruwa har ma manyan jiragen ruwa sukan kafe wani lokaci.
Duk da shirin taimako, da saurin isowar \VUgras{jiragen ceto,}
\VUgras{nutsewar jiragen ruwan} nan tana da ban tsoro, kuma, a lokacin
\VUgras{guguwa} ta daidaita rana da dare, jiragen ruwa da yawa sun hallaka
tare da mutane da kaya.

%%K t10-25-2 p
Duk da wannan makwabtaka, \VUgras{tashar jiragen ruwa} tana da yawan
matuƙan jirgin ruwa ƙwarai.

\VUsaut{6.0mm}
\VUfilet{22mm}
